\chapter*{Reader's Guide}
\addcontentsline{toc}{chapter}{Reader's Guide}

This book develops a single framework across multiple levels of abstraction, from formal mathematical results to applications in cognition, institutions, and cosmology. Readers interested in different aspects of the project may wish to approach the material through different paths.

\section*{The Structure of the Book}

The book is organized into five parts.

\textbf{Part One} develops the mathematical foundation. It introduces trajectory spaces, compression operators, interaction decompositions, stability sequences, reconstruction operators, and the Fundamental Error Decomposition. Readers interested primarily in the formal structure of the framework should begin here. The four chapters of Part One are self-contained and do not presuppose familiarity with any of the application domains.

\textbf{Part Two} instantiates the framework in three domains where the conditions for reconstruction can be established rigorously: learning systems, semantic inference, and physical field systems. These chapters develop the framework's verified classes, apply the five diagnostic questions to each domain, and identify the open empirical obligations that each domain generates. Chapter~5 (learning systems) is the most accessible entry point to Part Two; Chapter~7 (physical field systems) is the most technically demanding.

\textbf{Part Three} extracts a domain-independent diagnostic methodology from the results of Parts One and Two. The five diagnostic questions, the failure taxonomy, the alignment design principle, and the formal audit procedure provide the practical tools used throughout the remainder of the book.

\textbf{Part Four} applies the framework to domains that do not belong to the verified classes of Part Two: individual memory, institutional records, and cosmology. These chapters demonstrate the framework's research-program-generating function. The conclusions are frequently conditional, and open empirical obligations appear throughout. This is not a weakness---it is the framework operating correctly in domains where the evidence is incomplete.

\textbf{Part Five} develops the philosophical synthesis. It examines why the past matters less at higher interaction orders, what projection determines about which distinctions survive compression, and what the geometry of alternative histories reveals about the conditions for knowledge, explanation, and memory.

\section*{Recommended Reading Paths}

\textbf{For readers interested in the mathematics and formal theory:} Read Parts One, Two, and Five in sequence.

\textbf{For readers interested in machine learning, semantic systems, or physical field modeling:} Read Part One in full, then focus on the relevant chapter in Part Two. Chapters 8 and 9 in Part Three develop the failure analysis and alignment principles that extend the Part Two results.

\textbf{For readers interested in cognition or memory:} Read Chapter~1, Chapter~10, and Chapter~11. Chapters 2--4 and Part Five provide formal and philosophical context.

\textbf{For readers interested in institutions or historical records:} Follow the same path as above, substituting Chapter~12. The failure-mode analysis of Chapter~8 maps particularly directly onto institutional phenomena.

\textbf{For readers interested in cosmology or philosophy of physics:} Read Chapter~1, Chapter~4, Chapter~7, and Chapter~13. Chapter~14 connects stability to the philosophy of explanation.

\textbf{For readers interested primarily in the framework as a methodology:} Begin with Chapter~10, which provides the complete five-question audit procedure, then use earlier chapters as references when formal justification is needed.

\section*{Three Distinctions to Keep in Mind}

A \textbf{verified result} is a theorem, proof, or empirically supported claim established within a specified class of systems. Verified results appear primarily in Parts One and Two.

A \textbf{conditional result} is a statement whose validity depends on assumptions that remain empirically or mathematically open. Conditional results appear throughout Parts Three, Four, and Five.

An \textbf{open empirical obligation} is a measurement, simulation, proof, or experiment whose completion would resolve a conditional result. Open obligations are outputs of the framework: precise scientific questions generated by applying it to a domain where evidence is currently insufficient.

\section*{How to Use This Book}

The chapters that follow move from mathematics to application and from application to philosophy. The order is deliberate. Each layer provides concepts that become precise only after the previous layer has been developed.

Readers interested in the destination may begin almost anywhere.

Readers interested in the foundations should begin at the beginning.
